\chapter{Hướng phát triển}
\label{ch:huong-phat-trien}

\section{Cải tiến sản phẩm}
\label{sec:huong-phat-trien}

\paragraph{Đã xong trong bản này} (so với danh sách ``hướng phát triển'' của bản báo cáo trước):
tách thư mục chỉ mục theo profile; fp16 cho embedding và reranker; model server HTTP cho
embedding/reranker; mở rộng lên 6 văn bản; giao diện Gradio; ghim phiên bản thư viện; kiểm thử tự
động; đánh giá có khoảng tin cậy, baseline và kiểm định cặp; câu trả lời tất định
(\texttt{temperature=0}, \texttt{seed}). Trong tuần 05--09/10: chạy LLM và model server trên Colab qua
Cloudflare tunnel; ghi nhật ký câu hỏi của giao diện Gradio; chạy lại 56 câu với prompt mới; thêm chú
giải thang điểm vào prompt và công tắc \texttt{RAG\_THINKING}.

\paragraph{Còn lại, theo thứ tự ưu tiên:}
\begin{enumerate}
  \item \textbf{Chạy} \texttt{make eval-answers} (LLM trên Colab) để có số faithfulness/relevancy/
        correctness, và so chế độ RAG với \texttt{closed-book}. Đây là việc còn thiếu duy nhất của
        phần đánh giá. Cùng đợt chạy đó, chạy lại 56 câu theo từng yếu tố một (chỉ đổi prompt, chỉ đổi
        top-$k$ đưa vào LLM, bật/tắt \texttt{RAG\_THINKING}) để tách tác dụng của từng thay đổi ở
        Mục~\ref{sec:chay-lai-56-cau}.
  \item \textbf{Xử lý khóa tuyển sinh}: gắn metadata phạm vi áp dụng (ví dụ ``khóa 2025 trở đi'') cho
        mục sửa đổi của QĐ 3183, và yêu cầu câu trả lời về điểm nêu rõ áp dụng cho khóa nào (hoặc nêu
        cả hai bảng). Cùng lúc, gắn câu dẫn ``thang điểm 10''/``quy đổi để tính điểm trung bình'' vào
        mục sửa đổi để thay chú giải đang đặt cứng trong prompt.
  \item \textbf{Mở rộng bộ câu hỏi} lên vài trăm câu, có nhãn loại câu (trong/ngoài phạm vi, viết tắt,
        nối tiếp, cần tra bảng) và hai người gán nhãn độc lập.
  \item \textbf{Điền metadata hiệu lực} (\texttt{so\_hieu}, \texttt{ngay\_ban\_hanh}, \texttt{tinh\_trang})
        rồi dùng \texttt{MetadataFilters} để loại văn bản hết hiệu lực --- \texttt{processed\_loader}
        đã sẵn sàng nhận các trường này từ front matter.
  \item \textbf{Liên kết Điều gốc với mục sửa đổi} bằng \texttt{NodeRelationship}: khi truy hồi được
        Điều 7 của Quy chế 2021 thì tự động kèm mục ``sửa đổi Điều 7'' của QĐ 3183, và ngược lại.
  \item \textbf{Small-to-big}: cắt thêm Node con theo khoản, truy hồi trên Node con rồi đưa Điều cha
        cho LLM (\texttt{HierarchicalNodeParser} + \texttt{AutoMergingRetriever}) --- nhắm đúng lỗi
        ``đọc sai chi tiết số'' trong Điều dài.
  \item \textbf{Chat engine có ngữ cảnh} (\texttt{CondensePlusContextChatEngine}) cho câu hỏi nối tiếp;
        giao diện Gradio đã có sẵn lịch sử chat.
  \item \textbf{Công cụ tra bảng quy đổi điểm} (\texttt{FunctionTool} + agent) thay vì để LLM tự tính.
  \item \textbf{Vector store ngoài} (Chroma/Qdrant) khi corpus vượt vài trăm Node, để lọc metadata
        phức tạp và cập nhật từng phần.
\end{enumerate}


\section{Định hướng bản thân}

\chualam{từng thành viên viết ở bản báo cáo cuối}.

\section{Kết luận}

Trên corpus 6 văn bản / 54 Node, hệ thống đạt Recall@1 = 81\% và Recall@3 = 100\% ở bộ 31 câu hỏi về
Quy chế 2021, và Recall@1 = 90\% ở bộ 10 câu về nội dung sửa đổi --- trong khi baseline tìm từ khóa
BM25 chỉ đạt 58\% và 40\% ở R@1. Chỉ số do hệ thống tự tính khớp với chỉ số do
\texttt{RetrieverEvaluator} của LlamaIndex tính. Hai kỹ thuật ``nâng cao'' thường được khuyến nghị
--- xếp hạng lại bằng cross-encoder và hybrid BM25 + vector --- \emph{không} cải thiện trên quy mô
này và có thể làm giảm chất lượng nếu chọn sai tham số trộn, nên cả hai đều mặc định tắt và có công
tắc đo lại bằng \texttt{make eval}. Điểm yếu hiện nay không nằm ở truy hồi vector mà ở ba chỗ: dữ
liệu (54 Node, metadata hiệu lực trống), tầng sinh (LLM đọc sai số liệu, không nhớ hội thoại), và
độ mạnh thống kê của bộ câu hỏi đánh giá.
